% =====================================================================
%  CAPÍTULO 1 -- INTRODUÇÃO
%  Formato: PROJETO APLICADO (relatório técnico + protótipo)
%
%  RASCUNHO INICIAL redigido a partir do pré-projeto (v3). Revise o texto,
%  ajuste ao que de fato foi feito e remova as caixas de orientação ao final.
%
%  O Manual recomenda texto corrido, sem dividir a introdução em seções.
% =====================================================================

\chapter{Introdução}
\label{cap:introducao}

% (a) Contexto e problema
A prática regular de lutas e artes marciais está associada a benefícios
que vão além do condicionamento físico. Em revisão sistemática com
meta-análise, \citeonline{moore2020} encontraram efeitos positivos do
treinamento em artes marciais sobre o bem-estar e a saúde mental dos
praticantes, incluindo a redução de sintomas de ansiedade e depressão.
Esses benefícios, porém, só aparecem com a continuidade do treino. Na
prática das academias, uma parte expressiva dos alunos abandona a
atividade antes de consolidar o hábito, seja pela percepção lenta da
própria evolução, pela falta de metas claras, pelo distanciamento do
professor ou por mudanças na rotina. Para o aluno, o abandono interrompe
o ciclo em que os ganhos se tornariam perceptíveis; para a academia e o
professor, representa perda de receita e de vínculo com a turma.

% (b) Solução proposta
Para enfrentar esse problema, este projeto propõe o \textbf{RoundUp}, um
aplicativo móvel que aplica \textbf{gamificação} -- o uso de elementos de
design de jogos em contextos que não são jogos \cite{deterding2011} -- ao
acompanhamento dos treinos. Para o aluno, o aplicativo oferece mecânicas
como pontos de experiência, sequências de treino, medalhas e uma barra de
progresso rumo à próxima graduação. Para o professor, oferece apoio à
gestão: registro de presença, criação de missões para a turma e alertas
sobre alunos com queda de frequência. O sistema distingue três perfis de
usuário: o \textbf{aluno}, o \textbf{professor}, que pode atuar de forma
autônoma, e a \textbf{academia}, que reúne professores, turmas e alunos.
O protótipo resultante é de natureza \textbf{digital} (Regulamento,
Art.~10): um aplicativo funcional em React Native com \textit{back-end}
no Firebase.

% Natureza de ideação do projeto
Por se tratar de um domínio em que o valor de cada mecânica de jogo para
o público ainda é pouco conhecido, o projeto foi conduzido como um
\textbf{projeto de ideação}. Em vez de partir de requisitos fechados, o
autor elaborou um conjunto de ideias de funcionalidades e as converteu em
\textbf{requisitos iniciais}. Esses mesmos requisitos tiveram dois usos
simultâneos: orientaram a construção do aplicativo e se tornaram os itens
de um questionário de avaliação de ideias (\textit{survey}), respondido
por praticantes e ex-praticantes reais. Como houve poucos respondentes
(seis) e nenhum professor entre eles, foram construídas 12 personas
sintéticas e, a partir delas, geradas 65 \textbf{respostas sintéticas} ao
questionário, com apoio de um modelo de linguagem de grande porte (LLM).
Essas respostas foram usadas para testar o instrumento e o procedimento
de análise, e não como evidência da percepção de usuários. A \autoref{fig:ciclo-ideacao}, no
\autoref{cap:requisitos}, ilustra esse ciclo.

% (c) Objetivos
O objetivo geral deste projeto é \textbf{idear, desenvolver e avaliar o
RoundUp, um aplicativo gamificado para academias de lutas e artes
marciais, voltado a aumentar o engajamento e a retenção dos alunos e a
apoiar o professor na gestão da turma}. Os objetivos específicos são:

\begin{enumerate}
  \item identificar fatores de motivação e de evasão de praticantes de
        lutas e artes marciais;
  \item elaborar ideias de mecânicas de gamificação e de apoio ao
        professor e convertê-las em requisitos iniciais rastreáveis;
  \item avaliar a percepção de valor dessas ideias com usuários reais,
        por meio de um questionário;
  \item construir personas sintéticas e usá-las como insumo
        complementar da ideação e para gerar respostas de teste que
        verifiquem o questionário e o procedimento de análise, sem
        misturá-las aos dados reais;
  \item modelar e implementar um protótipo do aplicativo com os
        requisitos priorizados;
  \item considerar, desde a concepção, a conformidade com a Lei Geral de
        Proteção de Dados Pessoais (LGPD), inclusive no tratamento de dados
        de menores de idade \cite{lgpd2018}.
\end{enumerate}

% (d) Justificativa
O projeto tem relevância social, por contribuir para a permanência de
alunos em uma atividade associada à saúde física e mental; relevância
técnica, por articular engenharia de requisitos, gamificação,
desenvolvimento móvel e proteção de dados pessoais em uma solução
prática; e relevância metodológica, por documentar o uso de personas e
respostas sintéticas como insumo complementar da ideação e como dados de
teste de um instrumento de pesquisa, deixando explícito o que elas podem
e o que não podem mostrar. Soma-se a isso a vivência do autor
com as artes marciais, que facilita o acesso a praticantes e professores.

% (e) Escopo
O escopo do projeto concentra-se em academias de lutas e artes marciais
no Brasil e em duas perspectivas de uso: a do aluno (engajamento) e a do
professor (gestão). Ficam fora do escopo: outras modalidades esportivas,
prescrição de treinos ou orientação nutricional, integração com
dispositivos vestíveis, análise de movimentos por visão computacional e
aspectos comerciais, como monetização e escalabilidade em produção. A
avaliação do impacto de longo prazo sobre a evasão também ultrapassa o
prazo do TCC.

O pré-projeto citava Boxe e Muay Thai como foco inicial. O recorte foi
ampliado para lutas e artes marciais em geral: o aplicativo não depende de
uma modalidade específica, e o usuário escolhe a sua no cadastro.

% (f) Estrutura do relatório
Este relatório está organizado da seguinte forma: o
\autoref{cap:requisitos} apresenta os fundamentos que orientaram a
ideação, o processo de ideação e os requisitos do projeto; o
\autoref{cap:modelagem}, a modelagem da solução; o
\autoref{cap:implementacao}, as estratégias de implementação do
protótipo; o \autoref{cap:testes}, os procedimentos de avaliação das
ideias com usuários reais, o teste do questionário com respostas
sintéticas e os testes do protótipo; e o \autoref{cap:resultados}, os resultados e as
considerações finais.
